\subsection{背景与资源模型}

在同一空域、同一时段内工作的用频装备越来越多，若两部装备在重叠的频段上同时发射，就会互相干扰。
为在计划阶段（而非事后）消除干扰，需要把“谁在什么时间用哪一段频率”的全部申请汇总起来，
先\emph{检测}其中潜在的时频冲突，再通过调整少量计划参数或撤销个别计划把冲突\emph{消解}掉。

题目将时频资源离散化：时域以 $\Delta t$ 为最小不可拆分单位，频域把有限带宽等分为 $100$ 个宽度为 $\Delta f$ 的频段，
编号 $0,1,\dots,99$。于是“时频资源”就是一张以频段编号为行、以时刻为列的二维栅格，
每个栅格单元 $(b,t)$ 表示“第 $b$ 个频段在第 $t$ 个时隙”这一份最小资源。
冲突即两个用频计划占用了同一个单元。

\subsection{用频计划的参数与冲突的判定}

附件 1 给出 $\valQonePlans$ 个用频装备提报的用频计划，按编号首字母分为 A、B、C 三类
（各类装备数与参数见表~\ref{tab:data}）。每个计划由四个参数确定：

\begin{enumerate}[label=(\roman*),leftmargin=2.4em,itemsep=0.15em]
  \item \textbf{频段区间}：连续占用的频段 $[f,\,f+w)$，占用的频段数 $w$ 在任何调整下保持不变；
  \item \textbf{首次使用的时间区间}：$[s,\,s+d)$，$d$ 为每次使用的时长；
  \item \textbf{间隔时长} $g$：相邻两次使用之间的空闲时长；
  \item \textbf{使用次数} $n$：按上述规则重复使用 $n$ 次。
\end{enumerate}

题目以装备 A001 为例规定了推进规则：首次使用 $[35,40)$、间隔 $60$，则第二次使用为 $[40+60,\,45+60)=[100,105)$。
即第 $k$ 次（$k=0,1,\dots,n-1$）使用的时间区间为 $[\,s+k(d+g),\ s+d+k(d+g)\,)$，
所有区间均为半开区间。\textbf{时频冲突}的判定是：两个用频计划的频段区间存在交叠，
\emph{且}它们的某一次使用在时间上存在交叠。

\subsection{消解规则}

题目对消解给出了如下规则，本文将其区分为\emph{硬约束}与\emph{优先目标}两类：

\begin{itemize}[leftmargin=1.6em,itemsep=0.15em]
  \item \textbf{硬约束}：一个用频计划只能调整其中一个参数或撤销；一般不调整使用次数和使用间隔时长
        （问题 4 对 C 类装备的间隔时长解除该限制）；频段、时间调整的最大平移幅度分别为 $10\Delta f$、$5\Delta t$；
        消解后的全部计划之间不得存在任何时频冲突。
  \item \textbf{优先目标}（按题目给出的先后次序）：尽可能减少撤销的用频计划；尽量减少被调整的用频计划数量；
        A 类装备优先级最高、B 类其次、C 类最低，高优先级装备的计划应尽量保持；尽量减少调整的幅度。
\end{itemize}

\subsection{四个问题与交付物}

\begin{description}[leftmargin=0em,itemsep=0.2em,font=\normalfont\bfseries]
  \item[问题 1（检测）] 对附件 1 的全部计划进行时频冲突检测，把每一对冲突写入 \nolinkurl{result1.xlsx}，
        并在论文中给出冲突的统计结果。
  \item[问题 2（消解）] 对检测出的冲突进行消解，把需要调整或撤销的装备写入 \nolinkurl{result2.xlsx}
        （B 列填调整后频段区间、C 列填调整后时间区间、D 列填“是”表示撤销），
        并按题目表 1 的格式给出各类别的保留 / 调整 / 撤销数量。
  \item[问题 3（增容）] 基于问题 2 得到的无冲突方案，在不限制频段和时间平移幅度、不增加时频资源的前提下，
        最多还能安排多少个 C 类用频装备？给出相应计划并写入 \nolinkurl{result3.xlsx}。
  \item[问题 4（放宽间隔）] 允许部分 C 类装备调整用频间隔时长（与原间隔的差异不超过 $10\Delta t$），
        重新消解问题 1 检测出的冲突，写入 \nolinkurl{result4.xlsx} 并给出表 1 格式的统计结果。
\end{description}

问题之间是串行依赖的：问题 1 的冲突集合是问题 2、4 的输入，问题 2 的无冲突方案是问题 3 的输入。
本文把四个问题分别形式化为\emph{区间相交判定}、\emph{带类别优先级的词典序整数规划}、
\emph{二维集合装填}与\emph{动作集合扩张后的同一词典序规划}，并对每个结果给出与求解器分离的独立校验。
